\documentclass[11pt]{article}
\usepackage[margin=1in]{geometry}
\usepackage{amsmath, amssymb, amsthm}
\usepackage{hyperref}

\title{Fixed Income Mathematics:\\Curves, Bonds, Duration, and Everything In Between}
\author{Abhi Wadhwa}
\date{}

\begin{document}
\maketitle

\begin{abstract}
My notes on the math behind fixed income analytics: yield curve bootstrapping, bond pricing, duration and convexity, and the Nelson-Siegel-Svensson parametric model. I wrote these while building a fixed income toolkit and wanted to keep track of where the formulas come from and---more importantly---why they make sense. The goal is to go from raw market data (deposit rates, swap rates) all the way to portfolio risk metrics (DV01, key rate durations) with clear reasoning at each step.
\end{abstract}

\section{The Yield Curve: Starting From Scratch}

The yield curve is the foundation of everything in fixed income. It tells you the time value of money---how much a dollar received at time $t$ in the future is worth today. The fundamental object is the \textbf{discount factor} $D(t)$: the present value of \$1 to be received at time $t$.

But we don't observe $D(t)$ directly. We observe prices and rates of traded instruments, and we need to extract the discount curve from them. This process is called \textbf{bootstrapping}.

\section{Yield Curve Bootstrapping}

The idea is to build the discount curve sequentially, starting from the short end and working outward. At each step, we use a new market instrument to determine one more point on the curve.

\subsection{Deposits (Short End)}

For maturities up to about 1 year, we use deposit (money market) rates. A deposit at rate $r$ for maturity $T$ pays simple interest, so:
\[
D(T) = \frac{1}{1 + r \cdot T}
\]

This is the simplest instrument---you lend \$1, get back $1 + rT$, so the present value of that future payment is $1/(1+rT)$.

\subsection{Forward Rate Agreements (FRAs)}

FRAs lock in a rate for a future period $[T_1, T_2]$. If we already know $D(T_1)$ from the deposit curve, we can extract $D(T_2)$ using the no-arbitrage relationship:
\[
D(T_2) = \frac{D(T_1)}{1 + r_{\text{FRA}} \cdot (T_2 - T_1)}
\]

The logic: investing from now to $T_2$ should give the same result whether you go directly or go to $T_1$ first and then forward from $T_1$ to $T_2$.

\subsection{Interest Rate Swaps (Long End)}

This is where it gets interesting. A par swap has the property that its present value is zero at inception, which means the PV of the fixed leg equals the PV of the floating leg. For a swap with annual fixed payments at rate $c$ and maturity $T$:
\[
\sum_{i=1}^{n} c \cdot D(t_i) + D(T) = 1
\]

The left side is the PV of the fixed leg (coupons plus notional), and the right side is the PV of the floating leg (which equals par at inception---this is one of those facts that seems like magic the first time you see it but becomes obvious once you understand that floating rate bonds always price at par on reset dates).

If we already know $D(t_1), \ldots, D(t_{n-1})$ from shorter instruments, we can solve for $D(T)$:
\[
D(T) = \frac{1 - c \sum_{i=1}^{n-1} D(t_i)}{1 + c}
\]

I spent a while trying to figure out why the floating leg equals par, and the trick is this: on each reset date, the floating rate is set to make the next period's payment fair, so the bond is worth par right after each reset. Working backwards from maturity, this means it's worth par at every reset date, including today.

\subsection{Zero Rates and Interpolation}

Once we have discrete discount factors, we can extract the continuously compounded zero rate:
\[
z(t) = -\frac{\ln D(t)}{t} \qquad \Leftrightarrow \qquad D(t) = e^{-z(t) \cdot t}
\]

Between bootstrapped points, we interpolate $z(t)$. Cubic spline interpolation on zero rates is common and well-behaved. But for a smooth, parametric representation, we can use Nelson-Siegel-Svensson.

\section{Nelson-Siegel-Svensson Model}

The NSS model parameterizes the zero rate curve with six parameters $(\beta_0, \beta_1, \beta_2, \beta_3, \lambda_1, \lambda_2)$:
\[
z(\tau) = \beta_0 + \beta_1 \cdot \frac{1 - e^{-\tau/\lambda_1}}{\tau/\lambda_1} + \beta_2 \cdot \left[\frac{1 - e^{-\tau/\lambda_1}}{\tau/\lambda_1} - e^{-\tau/\lambda_1}\right] + \beta_3 \cdot \left[\frac{1 - e^{-\tau/\lambda_2}}{\tau/\lambda_2} - e^{-\tau/\lambda_2}\right]
\]

Each parameter has a clear interpretation, which is what makes this model so elegant:

$\beta_0$ is the \textbf{long-term level}---as $\tau \to \infty$, $z(\tau) \to \beta_0$. It controls where the curve asymptotes.

$\beta_1$ governs the \textbf{short-term component}. The instantaneous short rate is $\beta_0 + \beta_1$. If $\beta_1 < 0$, the curve slopes upward (normal); if $\beta_1 > 0$, it slopes downward (inverted).

$\beta_2$ creates a \textbf{hump or trough} at maturity $\lambda_1$. This captures the medium-term curvature that you see in most real yield curves.

$\beta_3$ adds a \textbf{second hump} at maturity $\lambda_2$, giving extra flexibility to fit more complex curve shapes.

Parameters are fitted by least squares:
\[
\min_{\theta} \sum_{i} w_i \cdot \big(z_{\text{model}}(t_i; \theta) - z_{\text{market}}(t_i)\big)^2
\]

The optimization is nonlinear in $\lambda_1$ and $\lambda_2$ (linear in the betas, conditional on the lambdas), so in practice people often do a grid search over the lambdas and solve for the betas analytically at each grid point.

\section{Bond Pricing}

A fixed-coupon bond with face value $F$, coupon rate $c$, $m$ payments per year, and $n$ total coupon periods is priced as:
\[
P = \sum_{i=1}^{n} \frac{c \cdot F / m}{(1 + y/m)^{i}} + \frac{F}{(1 + y/m)^{n}}
\]

where $y$ is the yield to maturity.

\textbf{Yield to maturity} is the single discount rate that equates the PV of all cash flows to the market price. It's the internal rate of return of the bond, and it has to be solved numerically (there's no closed form for $n > 4$). Brent's method works well here---it's bracketed, doesn't need derivatives, and converges quickly.

One thing worth noting: yield to maturity is a somewhat misleading number because it assumes all coupons are reinvested at rate $y$, which generally won't be the case. But it's the convention, and everyone uses it.

\section{Duration}

Duration measures how sensitive a bond's price is to changes in yield. There are several flavors, and they're all related.

\subsection{Macaulay Duration}

\textbf{Macaulay duration} is the weighted-average time to receive the bond's cash flows, where the weights are the present values of each cash flow:
\[
D_{\text{mac}} = \frac{1}{P} \sum_{i=1}^{n} t_i \cdot CF_i \cdot D(t_i)
\]

Intuitively, a bond with cash flows concentrated further in the future has higher duration and is more sensitive to rate changes. A zero-coupon bond has Macaulay duration equal to its maturity---there's only one cash flow, so the weighted average is trivially the maturity date.

\subsection{Modified Duration}

\textbf{Modified duration} converts Macaulay duration into a direct measure of price sensitivity:
\[
D_{\text{mod}} = \frac{D_{\text{mac}}}{1 + y/m}
\]

The key relationship is:
\[
\frac{dP}{P} \approx -D_{\text{mod}} \cdot dy
\]

So a bond with modified duration 5 loses approximately 5\% of its value for a 1\% (100 bps) increase in yield. This is the most commonly used risk metric in fixed income.

\subsection{DV01}

\textbf{DV01} (Dollar Value of a Basis Point) is the dollar change in price for a 1 bp move in yield:
\[
\text{DV01} = D_{\text{mod}} \cdot P \cdot 0.0001
\]

DV01 is more useful than duration for hedging because it's in dollar terms. If I'm long a bond with DV01 = \$50 and want to hedge with a bond that has DV01 = \$25, I need a 2:1 ratio.

\section{Convexity}

Duration gives a linear approximation to the price-yield relationship, but the actual relationship is convex (curved). \textbf{Convexity} captures this curvature:
\[
C = \frac{1}{P} \sum_{i=1}^{n} t_i \cdot (t_i + 1/m) \cdot \frac{CF_i}{(1 + y/m)^{t_i \cdot m + 2}}
\]

The improved price approximation using both duration and convexity is:
\[
\boxed{\Delta P \approx -D_{\text{mod}} \cdot P \cdot \Delta y + \frac{1}{2} \cdot C \cdot P \cdot (\Delta y)^2}
\]

The convexity term is always positive for vanilla bonds (whether rates go up or down, the second-order effect helps you). This means \textbf{convexity is always good to have}---for the same duration, a more convex bond outperforms in large rate moves.

After playing with this for a while in the toolkit, I found that the duration-only approximation is fine for moves up to about 50 bps, but for larger moves (100+ bps), the convexity correction makes a meaningful difference. For a 10-year bond with duration 8 and convexity 80, a 200 bp move gives a duration-only estimate of $-16\%$ but the convexity-adjusted estimate is $-16\% + 1.6\% = -14.4\%$. That 1.6\% difference is real money on a large portfolio.

\section{Putting It All Together}

The workflow from market data to portfolio risk is:
\begin{enumerate}
\item Collect market rates (deposits, FRAs, swaps).
\item Bootstrap the discount curve $D(t)$.
\item Optionally fit NSS for a smooth parametric representation.
\item Price bonds using the discount curve.
\item Compute risk metrics (duration, convexity, DV01).
\item Run scenario analysis (parallel shifts, twists, butterflies) to estimate portfolio PnL under stress.
\end{enumerate}

Each step builds on the previous one, and the math is internally consistent---the same discount factors used for pricing are used for risk, ensuring no arbitrage in the analytics.

\begin{thebibliography}{9}
\bibitem{tuckman2012} Tuckman, B., Serrat, A. (2012). \textit{Fixed Income Securities: Tools for Today's Markets}. 3rd edition, Wiley.

\bibitem{nelson1987} Nelson, C.R., Siegel, A.F. (1987). ``Parsimonious Modeling of Yield Curves.'' \textit{Journal of Business}, 60(4), 473--489.

\bibitem{svensson1994} Svensson, L.E.O. (1994). ``Estimating and Interpreting Forward Interest Rates: Sweden 1992--1994.'' \textit{NBER Working Paper No.\ 4871}.

\bibitem{fabozzi2007} Fabozzi, F.J. (2007). \textit{Fixed Income Analysis}. 2nd edition, Wiley.

\bibitem{hull2018} Hull, J.C. (2018). \textit{Options, Futures, and Other Derivatives}. 10th edition, Pearson.
\end{thebibliography}

\end{document}
